\begin{frame}[fragile]
  \frametitle{Recaudo del Resto de Impuestos como \% del PIB (2025)}

  \begin{table}
  \centering
%  \scriptsize
  \begin{tabular}{l c}
  \hline
  \textbf{Impuesto} & \textbf{\% del PIB} \\
  \hline
  GMF (4x1000)                                  & 0,8 \\
  Aranceles                                     & 0,3 \\
  Impuestos a los combustibles y al carbono     & 0,1 \\
  Impoconsumo                                   & 0,2 \\
  Régimen SIMPLE                                & 0,1 \\
  Patrimonio                                    & 0,1 \\
  Resto                                         & 0,2 \\
  \hline
  \textbf{Total}                                & \textbf{1,9} \\
  \hline
  \end{tabular}
  \end{table}

  \vspace{0.2cm}
  {\centering\scriptsize
  Fuente: Balance fiscal del GNC, Ministerio de Hacienda y cálculos propios.\par
  Resto incluye impuestos saludables y a los plásticos, retención en la fuente de inmuebles, normalización y otros.\par
  El total se calcula con cifras sin redondear.\par}
  \note{[Sources] Figuras.xlsx, hoja BGNC, W6 (2025) y W9:W26, consultada el 18 de septiembre de 2026. Valores en porcentaje del PIB, no fracciones. GMF: W17 = 0,79658383652749; aranceles: W25 = 0,34183744201066585; combustibles y carbono: W13 + W14 = 0,09977857535951763; consumo: W15 = 0,2255939069817974; SIMPLE: W19 = 0,08888410261127896; patrimonio: W20 = 0,06741130553297259. Resto: W16 + W18 + W21 + W22 + W26 = 0,24429421159476283, que incluye CREE (cero en 2025), normalización, retención en la fuente de inmuebles, impuestos saludables y a los plásticos y otros impuestos. Total: suma de los siete rubros = W9 - W11 - W12 - W24 = 1,8643833806184853. Se excluyen renta e IVA interno y externo. Cifras mostradas con un decimal; la suma de las cifras redondeadas es 1,8, mientras que el total sin redondear se aproxima a 1,9.}

\end{frame}
